% =============================================================================
% Derivation of the q(F) Variational Update
% Self-Contained Reference for Supervised Bayesian MF Implementation
%
% Author:  Andrew Walther
% Date:    March 2026
%
% PURPOSE:
%   Derives the CAVI update for the biological factor matrix F.  Unlike the
%   L update, F appears only in the genomics likelihood (not survival).
%   A remarkable algebraic property: tau_j CANCELS in the EBNM pseudo-
%   observation x_j, but does NOT cancel in the pseudo-noise s_j.
%
% COMPANION CODE: code/update_F.R
% PARENT DOCUMENT: derivations/MF_UpdateDerivations/MF_Derivations_UpdateAlgo_REVISED.tex
% =============================================================================

\documentclass[11pt]{article}
\usepackage[margin=1in]{geometry}
\usepackage{amsmath, amssymb, amsthm}
\usepackage{bm}
\usepackage{xcolor}
\usepackage{hyperref}
\usepackage{booktabs}
\usepackage{array}
\usepackage{cancel}
\usepackage{tcolorbox}
\tcbuselibrary{skins, breakable}

% ---- Custom commands --------------------------------------------------------
\newcommand{\lbar}[2]{\bar{l}_{#1#2}}
\newcommand{\fbar}[2]{\bar{f}_{#1#2}}
\newcommand{\lsec}[2]{\overline{l^2_{#1#2}}}
\newcommand{\fsec}[2]{\overline{f^2_{#1#2}}}
\newcommand{\Rk}[0]{R^{-k}}
\newcommand{\revised}[1]{{\color{red}\textbf{[REVISED]:} #1}}

% ---- Coloured boxes ---------------------------------------------------------
\newtcolorbox{correctionbox}{
  colback=red!5!white, colframe=red!50!black,
  title=Correction Note, fonttitle=\bfseries\small, breakable
}
\newtcolorbox{ebnmbox}{
  colback=blue!5!white, colframe=blue!40!black,
  title=EBNM Problem, fonttitle=\bfseries\small, breakable
}
\newtcolorbox{derivbox}{
  colback=gray!6!white, colframe=gray!50!black,
  title=Derivation Steps, fonttitle=\bfseries\small, breakable
}
\newtcolorbox{verifybox}{
  colback=green!5!white, colframe=green!50!black,
  title=Verification Check, fonttitle=\bfseries\small, breakable
}
\newtcolorbox{codebox}{
  colback=yellow!5!white, colframe=orange!50!black,
  title=Code Mapping (\texttt{update\_F.R}), fonttitle=\bfseries\small, breakable
}

% ---- Document ---------------------------------------------------------------
\title{\textbf{Derivation of the $q(F)$ Variational Update}\\[4pt]
       \large Supervised Bayesian Matrix Factorization\\[4pt]
       \normalsize Self-Contained Reference for \texttt{code/update\_F.R}}
\author{Andrew Walther}
\date{March 2026}

\begin{document}
\maketitle

\begin{abstract}
This document derives the coordinate-ascent variational inference (CAVI)
update for the biological factor matrix $F$ in the Supervised Bayesian Matrix
Factorization model.  Unlike the $L$ update (which combines genomics and
survival), $F$ appears \emph{only} in the genomics likelihood.  The update
is a $p$-observation EBNM problem: each feature $j$ provides one pseudo-
observation $x_j$ and pseudo-noise $s_j$.

A central algebraic property is derived in Section~\ref{sec:cancellation}:
the feature-specific noise precision $\tau_j$ \emph{cancels} in the EBNM
pseudo-observation $x_j = B_F[j]/A_F[j]$, but does \emph{not} cancel in
the pseudo-noise $s_j = 1/\sqrt{A_F[j]}$.  This means all features receive
the same point estimate (weighted OLS) but features with higher precision
$\tau_j$ are subject to less EBNM shrinkage.

This is the third of four modular derivation documents
(q$\beta \to$ qL $\to$ qF (this) $\to$ q$\tau$).
\end{abstract}

\tableofcontents
\newpage

% =============================================================================
\section{Setup and Notation}\label{sec:setup}
% =============================================================================

\subsection{Model Summary}

The genomics likelihood is:
\begin{equation}
  Y_{ij} = \sum_{k=1}^K l_{ik} f_{jk} + \varepsilon_{ij},
  \qquad \varepsilon_{ij} \sim \mathcal{N}(0,\,\tau_j^{-1})
  \label{eq:genomics}
\end{equation}

\textbf{$F$ does not appear in the Cox survival term.}  Therefore the entire
$q(f_{jk})$ update is driven by the genomics likelihood alone.

\subsection{Key Notation}

\begin{center}
\begin{tabular}{lll}
\toprule
\textbf{Symbol} & \textbf{Definition} & \textbf{R variable} \\
\midrule
$\lbar{i}{k}$ & $\mathbb{E}_q[l_{ik}]$ & \texttt{EL[i,k]} \\
$\lsec{i}{k}$ & $\mathbb{E}_q[l_{ik}^2]$ & \texttt{EL2[i,k]} \\
$\fbar{j}{k}$ & $\mathbb{E}_q[f_{jk}]$ & \texttt{EF[j,k]} \\
$\fsec{j}{k}$ & $\mathbb{E}_q[f_{jk}^2]$ & \texttt{EF2[j,k]} \\
$\tau_j$ & Feature noise precision & \texttt{Tau[j]} \\
$\Rk_{ij}$ & Partial residual (excl.\ factor $k$) & \texttt{R\_k[i,j]} \\
$\Sigma_{EL2_k}$ & $\sum_i\lsec{i}{k}$ (scalar sum over samples) & \texttt{sum\_EL2\_k} \\
\bottomrule
\end{tabular}
\end{center}

\subsection{Partial Residual Matrix $\Rk$}

\begin{equation}
  \Rk_{ij} = Y_{ij} - \sum_{k'\neq k}\lbar{i}{k'}\fbar{j}{k'}
  \label{eq:Rk}
\end{equation}

This quantity was defined and computed in the $L$ update; the same
\texttt{compute\_R\_k} function is reused here.

% =============================================================================
\section{CAVI Objective for $q(f_{jk})$}\label{sec:objective}
% =============================================================================

The ELBO terms involving $f_{jk}$ are:
\begin{equation}
  \mathcal{F}[q(f_{jk})] \supset
  \mathbb{E}_q\!\left[\log P(Y \mid L, F, \tau)\right]
  + \mathbb{E}_{q_{f_k}}\!\left[\log\tfrac{g_{f_k}}{q_{f_k}}\right]
  \label{eq:elbo}
\end{equation}

The genomics log-likelihood for feature $j$ and patient $i$:
\begin{equation}
  \log P(Y_{ij} \mid f_{jk}, \ldots)
  \propto -\frac{\tau_j}{2}\mathbb{E}_q\!\left[(\Rk_{ij} - f_{jk}\lbar{i}{k})^2\right]
  \label{eq:llik-ij}
\end{equation}

% =============================================================================
\section{Derivation of $A_F$ and $B_F$}\label{sec:derivation}
% =============================================================================

\begin{derivbox}
Expanding the expectation in \eqref{eq:llik-ij} and noting that
$\mathbb{E}_q[l_{ik}^2] = \lsec{i}{k}$:

\begin{align*}
  \mathbb{E}_q\!\left[(\Rk_{ij} - f_{jk}l_{ik})^2\right]
  &= (\Rk_{ij})^2 - 2\,f_{jk}\,\Rk_{ij}\,\mathbb{E}_q[l_{ik}]
     + f_{jk}^2\,\mathbb{E}_q[l_{ik}^2] \\
  &= (\Rk_{ij})^2 - 2\,f_{jk}\,\Rk_{ij}\,\lbar{i}{k}
     + f_{jk}^2\,\lsec{i}{k}
\end{align*}

Multiply by $-\tau_j/2$ and sum over all samples $i = 1,\ldots,n$:
\begin{align}
  \sum_i -\frac{\tau_j}{2}\mathbb{E}_q[\cdots]
  &= -\frac{f_{jk}^2}{2}
       \underbrace{\tau_j\sum_i\lsec{i}{k}}_{A_F[j]}
     + f_{jk}
       \underbrace{\tau_j\sum_i\Rk_{ij}\lbar{i}{k}}_{B_F[j]}
     + \text{const}
  \label{eq:AF-BF}
\end{align}
\end{derivbox}

Therefore:
\begin{align}
  A_F[j] &= \tau_j \sum_i \lsec{i}{k} = \tau_j\,\Sigma_{EL2_k}
  \label{eq:AF} \\[4pt]
  B_F[j] &= \tau_j \sum_i \Rk_{ij}\lbar{i}{k}
           = \tau_j\,\bigl(R_k^\top\,\bar{\bm{l}}_k\bigr)_j
  \label{eq:BF}
\end{align}

where $\Sigma_{EL2_k} = \sum_i\lsec{i}{k}$ is a \emph{scalar}.

\begin{codebox}
\begin{verbatim}
sum_EL2_k <- sum(EL2_k)                         # scalar
A_F  <- pmax(Tau * sum_EL2_k, A_floor)          # p-vector [floored]
B_F  <- Tau * as.vector(t(R_k) %*% EL_k)        # p-vector
\end{verbatim}
\end{codebox}

% =============================================================================
\section{The $\tau$ Cancellation Property}\label{sec:cancellation}
% =============================================================================

This is a key structural property of the $F$ update.

\subsection{Cancellation in $x_j$}

The EBNM pseudo-observation is:
\begin{align}
  x_j &= \frac{B_F[j]}{A_F[j]}
       = \frac{\tau_j\,\bigl(R_k^\top\bar{\bm{l}}_k\bigr)_j}
              {\tau_j\,\Sigma_{EL2_k}}
  \label{eq:xj-cancel}
\end{align}

\begin{equation}
  \boxed{x_j = \frac{\bigl(R_k^\top\bar{\bm{l}}_k\bigr)_j}{\Sigma_{EL2_k}}
  \quad\text{--- \emph{independent} of }\tau_j}
  \label{eq:xj-final}
\end{equation}

\textbf{Interpretation:} $x_j$ is the weighted-least-squares estimate
of $f_{jk}$ that one would obtain if $l_{ik}$ were known exactly (i.e.,
in the OLS limit where $\lsec{i}{k} = \lbar{i}{k}^2$).  The noise level
$\tau_j$ does not affect the estimate itself.

\subsection{No Cancellation in $s_j$}

The EBNM pseudo-noise is:
\begin{equation}
  s_j = \frac{1}{\sqrt{A_F[j]}} = \frac{1}{\sqrt{\tau_j\,\Sigma_{EL2_k}}}
  \label{eq:sj}
\end{equation}

$\tau_j$ does \emph{not} cancel: features with higher precision $\tau_j$
(lower noise) receive smaller pseudo-noise $s_j$ and hence \emph{less}
EBNM shrinkage.

\begin{verifybox}
\textbf{Empirical check.}  With fixed $R_k$ and $EL_k$, compute
\texttt{update\_F\_k} twice:\ once with \texttt{Tau = rep(1, p)} and once
with \texttt{Tau = rep(1000, p)}.

\begin{itemize}
  \item \texttt{max|x(Tau=1) - x(Tau=1000)|} should be $\approx 0$ (machine epsilon).
  \item \texttt{mean(s(Tau=1)) / mean(s(Tau=1000))} should be $\approx \sqrt{1000} \approx 31.6$.
  \item Posterior means should differ (less shrinkage at $\tau=1000$).
\end{itemize}
This is verified in test T1.4 and T1.5 of \texttt{tests/test\_update\_F.R}
and demonstrated in Demo~1 of \texttt{demos/demo\_update\_F.R}.
\end{verifybox}

\begin{correctionbox}
\textbf{Correction R4} (from REVISED.tex): The EBNM subscript for the $F$
update should be $j$ (feature-indexed), not $i$.  The update produces a
$p$-vector of posterior means $\bar{f}_{jk}$ for $j = 1,\ldots,p$.
This was a typographical error in the original derivation document.
\end{correctionbox}

% =============================================================================
\section{Error-in-Variables: Loading Uncertainty in $A_F$}\label{sec:eiv}
% =============================================================================

The sum $\Sigma_{EL2_k} = \sum_i\lsec{i}{k}$ uses the \emph{full} second
moment of each loading, not the squared posterior mean.  This is the
error-in-variables correction for uncertain loadings:

\begin{equation}
  \sum_i\lsec{i}{k}
  = \sum_i\bigl(\mathrm{Var}_q(l_{ik}) + \lbar{i}{k}^2\bigr)
  = \underbrace{\sum_i\mathrm{Var}_q(l_{ik})}_{\geq 0}
    + \sum_i\lbar{i}{k}^2
  \geq \sum_i\lbar{i}{k}^2
  \label{eq:second-moment}
\end{equation}

Using $\lbar{i}{k}^2$ (pretending zero loading variance) would
underestimate $A_F$, overestimating the signal-to-noise ratio and
causing less shrinkage than appropriate.

% =============================================================================
\section{EBNM Formulation}\label{sec:ebnm}
% =============================================================================

\begin{ebnmbox}
\textbf{EBNM problem for $q(f_{\cdot k})$:}
\begin{equation}
  \mathrm{EBNM}\bigl(x_j = B_F[j]/A_F[j],\; s_j = 1/\sqrt{A_F[j]},\;
    g = g_{f_k}\bigr)
  \quad j = 1,\ldots,p
  \label{eq:ebnm}
\end{equation}
Calling convention:\quad
\texttt{ebnm(x = B\_F/A\_F, s = 1/sqrt(A\_F), prior\_family = "point\_normal")}\\[4pt]
Returns:\quad \texttt{posterior\$mean} $\to$ $\bar{f}_{jk}$,\quad
\texttt{posterior\$sd} $\to$ $\sigma_{jk}$\\
Second moment:\quad $\overline{f^2_{jk}} = \sigma_{jk}^2 + \bar{f}_{jk}^2$
\end{ebnmbox}

\begin{codebox}
\begin{verbatim}
x_F <- B_F / A_F                                # p-vector (tau-free!)
s_F <- 1.0 / sqrt(A_F)                          # p-vector (tau-dependent)
res <- ebnm(x = x_F, s = s_F, prior_family = prior_family)
f_mean   <- res$posterior$mean                  # p-vector
f_sd     <- res$posterior$sd                    # p-vector
f_second <- f_sd^2 + f_mean^2                   # p-vector
\end{verbatim}
\end{codebox}

% =============================================================================
\section{Gauss-Seidel Ordering in \texttt{update\_F\_all}}\label{sec:gauss-seidel}
% =============================================================================

The wrapper \texttt{update\_F\_all} iterates over $k = 1,\ldots,K$.
After updating $EF[\,,k]$, the new values are used when computing $\Rk$
for subsequent factors.  Importantly, the \emph{already-updated}
$EL[\,,k]$ from the preceding $L$ step is used here (Gauss-Seidel
across the L-then-F sequence).

% =============================================================================
\section{Verification Checks}\label{sec:verification}
% =============================================================================

\begin{verifybox}
\textbf{V1: OLS limit.}  When $\lsec{i}{k} = \lbar{i}{k}^2$ (zero loading
variance), $x_j = (R_k^\top\bar{\bm{l}}_k)_j / \sum_i\lbar{i}{k}^2$, which
is the OLS estimate of $f_{jk}$.

\textbf{V2: Null loading.}  When $\bar{l}_{ik} = 0$ and $\lsec{i}{k} = 0$
for all $i$, then $\Sigma_{EL2_k} = 0$, so $A_F$ hits the floor and $B_F = 0$,
yielding $\bar{f}_{jk} = 0$.

\textbf{V3: Scalar $\Sigma_{EL2_k}$.}  The loading second moments are summed
over all $n$ samples, giving a single scalar.  Therefore $A_F[j] = \tau_j
\times \text{const}$ — the only feature-variation in $A_F$ comes from $\tau_j$.

\textbf{V4: Second moment non-negativity.}  $\overline{f^2_{jk}} \geq
\bar{f}_{jk}^2$ element-wise, with equality iff posterior variance is zero.
\end{verifybox}

% =============================================================================
\section{Code Mapping}\label{sec:code}
% =============================================================================

\begin{center}
\begin{tabular}{lll}
\toprule
Equation & Code & Type \\
\midrule
$\Sigma_{EL2_k} = \sum_i\lsec{i}{k}$ & \texttt{sum(EL2\_k)} & scalar \\
$A_F[j] = \tau_j\,\Sigma_{EL2_k}$ & \texttt{Tau * sum\_EL2\_k} & $p$-vector \\
$B_F[j] = \tau_j(R_k^\top\bar{\bm{l}}_k)_j$ & \texttt{Tau * t(R\_k) \%*\% EL\_k} & $p$-vector \\
$x_j = B_F[j]/A_F[j]$ & \texttt{B\_F / A\_F} & $p$-vector ($\tau$-free) \\
$s_j = 1/\sqrt{A_F[j]}$ & \texttt{1/sqrt(A\_F)} & $p$-vector ($\tau$-dependent) \\
\bottomrule
\end{tabular}
\end{center}

% =============================================================================
\section{Algorithm Summary}\label{sec:algorithm}
% =============================================================================

\begin{enumerate}
  \item For $k = 1, \ldots, K$ (Gauss-Seidel order):
  \begin{enumerate}
    \item Compute partial residual:
      $\Rk = Y - EL\,EF^\top + \mathrm{outer}(EL[\,,k],\,EF[\,,k])$
    \item Compute loading second-moment sum (scalar):
      $\Sigma_{EL2_k} = \sum_i\lsec{i}{k}$
    \item Precision ($p$-vector):
      $A_F[j] = \max(\tau_j\,\Sigma_{EL2_k},\;\epsilon)$
    \item Signal ($p$-vector):
      $B_F[j] = \tau_j\,\bigl(R_k^\top\,\bar{\bm{l}}_k\bigr)_j$
    \item EBNM pseudo-obs ($\tau$-free): $x_j = B_F[j]/A_F[j]$
    \item EBNM pseudo-noise ($\tau$-dependent): $s_j = 1/\sqrt{A_F[j]}$
    \item Solve EBNM:
      $\mathrm{EBNM}(x_j, s_j)$ over $j = 1,\ldots,p$
    \item Update $\bar{f}_{jk} \leftarrow$ posterior mean,\quad
      $\overline{f^2_{jk}} \leftarrow \sigma_{jk}^2 + \bar{f}_{jk}^2$
  \end{enumerate}
\end{enumerate}

% =============================================================================
\appendix
\section{Notation Index}
% =============================================================================

\begin{tabular}{lll}
\toprule
Symbol & Meaning & Dimensions \\
\midrule
$Y$ & Genomics data & $n \times p$ \\
$L$ & Patient loadings & $n \times K$ \\
$F$ & Biological factors & $p \times K$ \\
$\tau_j$ & Feature noise precision & $p$-vector \\
$\Rk$ & Partial residual (excl.\ $k$) & $n \times p$ \\
$\Sigma_{EL2_k}$ & $\sum_i\lsec{i}{k}$ & scalar \\
$A_F[j]$ & EBNM precision for feature $j$ & $p$-vector \\
$B_F[j]$ & EBNM signal for feature $j$ & $p$-vector \\
$x_j$ & EBNM pseudo-obs ($\tau$-free) & $p$-vector \\
$s_j$ & EBNM pseudo-noise ($\tau$-dependent) & $p$-vector \\
\bottomrule
\end{tabular}

\end{document}
